\documentclass[tikz,border=4pt]{standalone}
\usepackage{resistencia}
\begin{document}
\begin{tikzpicture}[x=1cm,y=1cm]
% --- a) cuerpo en equilibrio con su plano de corte ---
\draw[cuerpo] (0.5,0) arc (270:180:0.5) -- (0,1.7) arc (180:90:0.5) -- (4.3,2.2) arc (90:0:0.5) -- (4.8,0.5) arc (0:-90:0.5) -- cycle;
\draw[fuerza] (-1.3,1.1) -- (-0.05,1.0) node[midway,above,etiq]{$\vec F_1$};
\draw[fuerza] (6.1,1.0) -- (4.85,1.1) node[midway,above,etiq]{$\vec F_2$};
\draw[fuerza] (1.2,-1.4) -- (1.2,-0.05) node[midway,right,etiq]{$\vec F_3$};
\draw[corte] (2.4,2.55) -- (2.4,-0.45);
\node[etiq,violeta,above] at (2.4,2.55) {plano de corte $\Pi$};
\node[etiq] at (2.4,-2.0) {a) Cuerpo en equilibrio};

% --- b) las dos partes separadas, con las fuerzas internas ---
\begin{scope}[xshift=8.2cm]
  % parte izquierda
  \draw[cuerpo] (2.4,0) -- (0.5,0) arc (270:180:0.5) -- (0,1.7) arc (180:90:0.5) -- (2.4,2.2) -- cycle;
  \draw[fuerza] (-1.3,1.1) -- (-0.05,1.0) node[midway,above,etiq]{$\vec F_1$};
  \draw[fuerza] (1.2,-1.4) -- (1.2,-0.05) node[midway,right,etiq]{$\vec F_3$};
  % la parte derecha empuja sobre la cara de corte (el cuerpo está comprimido por F1 y F2)
  \foreach \y/\l in {0.3/0.45, 0.8/0.7, 1.3/0.85, 1.85/0.5}
    \draw[vtension] ($(2.4,\y)+(\l,0)$) -- (2.4,\y);
  % parte derecha, separada
  \begin{scope}[xshift=4.2cm]
    \draw[cuerpo] (0,0) -- (1.9,0) arc (270:360:0.5) -- (2.4,1.7) arc (0:90:0.5) -- (0,2.2) -- cycle;
    \draw[fuerza] (3.7,1.0) -- (2.45,1.1) node[midway,above,etiq]{$\vec F_2$};
    \foreach \y/\l in {0.3/0.45, 0.8/0.7, 1.3/0.85, 1.85/0.5}
      \draw[vtension] ($(0,\y)+(-\l,0)$) -- (0,\y);
  \end{scope}
  \node[etiq,acero,align=center] at (3.8,3.15) {fuerzas internas:\\ cada parte actúa sobre la otra};
  \node[etiq] at (3.3,-2.0) {b) Dos partes, cada una en equilibrio};
\end{scope}
\end{tikzpicture}
\end{document}
